% 题证摘录；来源与整理边界见相同编号分题文件。
\paragraph{1. Branched Covers}
\begin{definition}[1]
For $(M,\omega)$ a symplectic manifold, $p\in M$, a local diffeomorphism $\phi:U\to\C^n$ for $U\ni p$ is $\omega$-tame if $(\phi_*\omega)(v,iv)>0$ $\forall v\ne0\in\C^n$. This is, $\phi^*J_0$ is $\omega$-tame, i.e. complex lines in $\C^n$ give symplectic submanifolds in $M$.
\end{definition}
\begin{definition}[2]
A map $f:X^4\to(Y^4,\omega_Y)$ from a compact, oriented manifold to a compact, symplectic manifold is a symplectic branched covering if $\forall p\in X$, $\exists U\ni p$, $V\ni f(p)$ coordinate neighborhoods (with $\phi:U\to\C^2$ an oriented diffeomorphism, $\psi:V\to\C^2$ an $\omega_Y$-tame diffeomorphism) s.t. the right vertical map of the commutative diagram
\[
\begin{tikzcd}
X\supset U\arrow[r,"\phi"]\arrow[d,"f"']&\C^2\arrow[d,"\psi f\phi^{-1}"]\\
Y\supset V\arrow[r,"\psi"']&\C^2
\end{tikzcd}\tag{1}
\]
is one of
\begin{enumerate}
\item $(u,v)\mapsto(z_1,z_2)=(u,v)$ (local diffeomorphism),
\item $(u,v)\mapsto(z_1,z_2)=(u^2,v)$ (simple branching),
\item $(u,v)\mapsto(z_1,z_2)=(u^3-uv,v)$ (cusp).
\end{enumerate}
\end{definition}
\begin{definition}[3]
The ramification curve is the set $R\subset X$ s.t. $R=\{p\in X\mid df(p)\text{ not onto}\}$. The branch (discriminant) curve is $D=f(R)\subset Y$, i.e. $D=\{q\in Y\mid\#f^{-1}(q)<\deg f\}$.
\end{definition}
We can calculate these curves explicitly in local coordinates. For instance, in the case of simple branching, we have that $\operatorname{Jac}(f)=\det(df)=\left|\begin{smallmatrix}2u&0\\0&1\end{smallmatrix}\right|=2u$, so $R=\{u=0\}$, $D=\{z_1=0\}$. In the case of a cusp, we have
\[
\operatorname{Jac}(f)=\det(df)=\begin{vmatrix}3u^2-v&-u\\0&1\end{vmatrix}=3u^2-v\tag{2}
\]
so $R=\{v=3u^2\}$ and $D=\{27z_1^2=4z_2^3\}$. What happens at the cusp: $\forall p\in R$, $\ker df=\C\times\{0\}\subset T_p\C^2$, so $\ker df$ is transverse to $TR$ at most points, but not at the cusp.
\begin{lemma}[1]
$R\subset X$ is a smooth, 2-dimensional submanifold, and $D\subset Y$ is a symplectic submanifold of $Y$ immersed except at isolated points (corresponding to cusps). In local coordinates, $D$ is a complex curve, so $TD$ consists of complex lines and $\omega_Y|_{TD}>0$.
\end{lemma}
\begin{proposition}[1]
If $f:X^4\to(Y^4,\omega_Y)$ is a symplectic branched covering, then $X$ carries a symplectic form $\omega_X$ (canonical up to isotopy) s.t. $[\omega_X]=f^*[\omega_Y]$.
\end{proposition}
\begin{proof}
Note that $f^*\omega_Y$ is a closed 2-form which is nondegenerate outside of $R$. $\forall p\in R$, $K_p=\ker df_p\subset T_pX$ is the kernel of $f^*\omega_Y$, and is a complex line in local coordinates (so it carries a natural orientation). We claim that $\exists\alpha$ an exact 2-form on $X$ s.t. $\forall p\in R$, $\alpha|_{K_p}>0$ is positive nondegenerate. Assuming this, we have that $\omega_X=f^*\omega_Y+\epsilon\alpha$ for $\epsilon>0$ sufficient small is closed and nondegenerate, since
\[
\omega_X\wedge\omega_X=f^*\omega_Y\wedge f^*\omega_Y+2\epsilon f^*\omega\wedge\alpha+\epsilon^2\alpha\wedge\alpha\tag{3}
\]
with the first term $\geq0$ everywhere and nondegenerate outside $R$, the second term positive on $R$ (in local coordinates, $f^*\omega_Y=\frac{i\lambda}{2}(dv\wedge d\bar v)$ for some $\lambda>0$, and $\alpha|_{\C\times\{0\}}>0$), and the third term negligible for small $\epsilon$.

We are left to prove our claim. Fix $p\in R$, and choose local coordinates $(u,v)$ on $X$ of our model. Set $x=\operatorname{Re}(u)$, $y=\operatorname{Im}(u)$, and $\alpha_p=d(\chi_1(|u|)\chi_2(|v|)xdy)$, where $\chi_1,\chi_2$ are cutoff functions chosen s.t. $\forall q\in R\cap\operatorname{Supp}(\chi_2)$, $\chi_1\equiv1$. In local coordinates, we have that
\[
\forall(u,v)\in R\cap\operatorname{Supp}(\chi_2),\qquad
\alpha_p|_{K=\C\times\{0\}}=\chi_2(|v|)dx\wedge dy>0\tag{4}
\]
and $0$ outside $\operatorname{Supp}(\chi_2)$. Covering $R$ by small neighborhoods, and taking $\alpha$ to be the sum of these $\alpha_p$ gives the desired exact form.

Finally, to see that the choice of $\omega_X$ is canonical up to isotopy, note that the set of $\alpha$'s satisfying our claim (i.e. exact 2-forms s.t. $\alpha|_K>0$ along $R$) is convex. That is, we can find an $\epsilon$ sufficiently small s.t., for two such forms $\alpha_1,\alpha_2$,
\[
f^*\omega_Y+\epsilon\alpha_1,\quad f^*\omega_Y+\epsilon\alpha_2,\quad
f^*\omega_Y+\epsilon(t\alpha_1+(1-t)\alpha_2))\tag{5}
\]
are all symplectic.
\end{proof}